\chapter{Аналитическая часть}
Целью лабораторной работы №3 является выделение участков программ, которые могут быть исполнены параллельно, на материале графовых моделей. Задание по варианту: найти максимальный элемент последовательности.

\section{Рекурсия}
\textbf{Рекурсия} --- это такой способ организации вычислений, при котором процедура или функция в ходе выполнения обращается сама к себе. Другими словами, рекурсия является методом определения функций (процедур), при котором определяемая функция применена в теле своего же собственного определения~\cite{lit1}.

\textbf{Рекурсивный алгоритм} --- алгоритм, определяемый через себя~\cite{lit2}.

\section{Граф}
Граф $G$ задаётся множеством точек или вершин $X=\{x_1, x_2, ..., x_n\}$ и
множеством линий $U=\{u_1, u_2, ..., u_m\}$, соединяющих между собой все или
часть этих точек. Таким образом, граф $G$ полностью задаётся парой $G=(X, U)$.

Если ребра из множества $U$ ориентированы, что обычно показывается
стрелкой $\vec{U}$,
то они называются дугами, и граф с такими ребрами
называется ориентированным графом, который обозначается $\vec{G} = (X,\vec{U})$. Если ребра не имеют ориентации, то такой граф называется неориентированным: $G=(X, U)$~\cite{lit3}.

\section{Графовые модели}
В графовых моделях программы представляются с помощью графов: набор вершин и множество соединяющих направленных дуг.

Вершины графа соответствуют программным элементам: процедурам, циклам, операторам, срабатывания операторов~\cite{lit4}.

Дуги отражают зависимости между этими элементами. Выделяют два основных типа связей:
\begin{itemize}
\item операционное отношение: дуга из вершины A в вершину B означает, что B может выполняться непосредственно после A;
\item информационное отношение: дуга из A в B указывает, что B использует данные, полученные в результате выполнения A.
\end{itemize}

 \section*{Вывод}
В аналитической части рассмотрены основные понятия связанные с графами и рекурсией.
\clearpage
